\documentclass[12pt]{article}

\usepackage[T1]{fontenc}
\usepackage[utf8]{inputenc}
\usepackage[romanian]{babel}
\usepackage{amsmath, amssymb, amsthm}
\usepackage{geometry}
\geometry{a4paper, margin=2.5cm}

\title{Grafuri și Arbori — Noțiuni Fundamentale și Algoritmi}
\author{}
\date{}

\begin{document}

\maketitle

\section{Introducere}

Grafurile sunt structuri fundamentale în informatică, folosite pentru a modela
rețele, drumuri, relații, dependențe, structuri ierarhice și multe altele.
Arborii sunt grafuri speciale, fără cicluri, care apar în mod natural în
structuri de date, algoritmi și modelarea ierarhiilor.

Acest capitol prezintă noțiunile de bază, tipurile de grafuri, reprezentări,
parcurgeri, arbori, arbori de căutare, arbori parțiali minimali și algoritmi
clasici precum BFS, DFS, Dijkstra, Kruskal și Prim.

\section{Grafuri — definiții fundamentale}

\subsection{Definiție}

Un \textbf{graf} este o pereche:


\[
G = (V, E),
\]


unde:
\begin{itemize}
    \item $V$ este mulțimea vârfurilor (nodurilor),
    \item $E$ este mulțimea muchiilor.
\end{itemize}

\subsection{Grafuri orientate și neorientate}

\begin{itemize}
    \item Într-un \textbf{graf neorientat}, muchiile sunt perechi neordonate $\{u,v\}$.
    \item Într-un \textbf{graf orientat} (digraf), muchiile sunt perechi ordonate $(u,v)$.
\end{itemize}

\subsection{Gradele unui nod}

\begin{itemize}
    \item În graf neorientat: gradul = numărul de muchii incidente.
    \item În graf orientat:
    

\[
    \deg^+(v) = \text{gradul de ieșire}, \quad
    \deg^-(v) = \text{gradul de intrare}.
    \]


\end{itemize}

\subsection{Lanțuri, drumuri, cicluri}

\begin{itemize}
    \item \textbf{Lanț}: succesiune de noduri conectate prin muchii.
    \item \textbf{Drum}: lanț fără repetarea nodurilor.
    \item \textbf{Ciclu}: drum închis.
\end{itemize}

\subsection{Conexitate}

\begin{itemize}
    \item Graf neorientat: \textbf{conex} dacă există drum între orice două noduri.
    \item Graf orientat: \textbf{tare conex} dacă există drum orientat între orice două noduri.
\end{itemize}

\section{Reprezentări ale grafurilor}

\subsection{Listă de adiacență}

Pentru fiecare nod păstrăm lista vecinilor săi.

Avantaje:
\begin{itemize}
    \item eficient pentru grafuri rare;
    \item spațiu $O(n + m)$.
\end{itemize}

\subsection{Matrice de adiacență}

Matrice $n \times n$ unde:


\[
A[i][j] = 
\begin{cases}
1, & \text{dacă există muchie } i \to j,\\
0, & \text{altfel}.
\end{cases}
\]



Avantaje:
\begin{itemize}
    \item verificare instantanee a existenței unei muchii;
    \item spațiu $O(n^2)$.
\end{itemize}

\section{Parcurgeri fundamentale}

\subsection{Parcurgerea în lățime (BFS)}

BFS explorează graful pe niveluri, folosind o coadă.

\subsubsection*{Pseudocod}

\begin{verbatim}
BFS(G, s):
  for each v in V:
    viz[v] = false
  queue = empty
  viz[s] = true
  enqueue(queue, s)

  while queue not empty:
    u = dequeue(queue)
    for each v in adj[u]:
      if not viz[v]:
        viz[v] = true
        enqueue(queue, v)
\end{verbatim}

\subsubsection*{Complexitate}



\[
O(n + m).
\]



\subsection{Parcurgerea în adâncime (DFS)}

DFS explorează cât mai adânc posibil înainte de a reveni.

\subsubsection*{Pseudocod}

\begin{verbatim}
DFS(u):
  viz[u] = true
  for each v in adj[u]:
    if not viz[v]:
      DFS(v)
\end{verbatim}

\subsubsection*{Complexitate}



\[
O(n + m).
\]



\section{Arbori}

\subsection{Definiție}

Un \textbf{arbore} este un graf neorientat, conex, fără cicluri.

Echivalent:
\begin{itemize}
    \item între orice două noduri există un drum unic;
    \item are $n$ noduri și $n-1$ muchii.
\end{itemize}

\subsection{Arbori înrădăcinați}

Un nod este ales ca rădăcină. Se definesc:
\begin{itemize}
    \item părintele,
    \item copiii,
    \item nivelul,
    \item înălțimea.
\end{itemize}

\section{Arbori binari}

\subsection{Definiție}

Un \textbf{arbore binar} este un arbore în care fiecare nod are cel mult doi copii:
\begin{itemize}
    \item copil stâng,
    \item copil drept.
\end{itemize}

\subsection{Tipuri de arbori binari}

\begin{itemize}
    \item complet,
    \item perfect,
    \item echilibrat,
    \item degenerat (listă).
\end{itemize}

\section{Arbori binari de căutare (BST)}

\subsection{Definiție}

Un arbore binar este BST dacă pentru orice nod $x$:


\[
\text{toate valorile din subarborele stâng} < x < \text{toate valorile din subarborele drept}.
\]



\subsection{Operații fundamentale}

\begin{itemize}
    \item căutare,
    \item inserare,
    \item ștergere.
\end{itemize}

\subsection{Complexități}



\[
\text{Best: } O(\log n),\quad
\text{Worst: } O(n).
\]



\section{Arbori parțiali minimali (MST)}

\subsection{Definiție}

Dat un graf ponderat neorientat, un \textbf{arbore parțial minimal} este un arbore
care conectează toate nodurile cu cost total minim.

\subsection{Algoritmul lui Kruskal}

Sortează muchiile crescător și le adaugă dacă nu formează ciclu (folosind Union-Find).

\subsubsection*{Complexitate}



\[
O(m \log m).
\]



\subsection{Algoritmul lui Prim}

Construiește MST extinzând treptat un arbore pornind dintr-un nod.

\subsubsection*{Complexitate}



\[
O(m \log n).
\]



\section{Drumuri minime}

\subsection{Algoritmul lui Dijkstra}

Calculează drumuri minime dintr-o sursă în grafuri cu ponderi pozitive.

\subsubsection*{Complexitate}



\[
O(m \log n).
\]



\subsection{Algoritmul Bellman-Ford}

Funcționează și cu ponderi negative.



\[
O(nm).
\]



\section{Grafuri orientate aciclice (DAG)}

\subsection{Definiție}

Un graf orientat fără cicluri se numește DAG.

\subsection{Ordonare topologică}

O ordonare a nodurilor astfel încât toate muchiile merg de la stânga la dreapta.

\subsubsection*{Pseudocod (Kahn)}

\begin{verbatim}
TopSort(G):
  compute indegree[v] for all v
  queue = all v with indegree 0
  while queue not empty:
    u = dequeue(queue)
    output u
    for each v in adj[u]:
      indegree[v]--
      if indegree[v] == 0:
        enqueue(queue, v)
\end{verbatim}

\section{Concluzii}

Grafurile și arborii sunt structuri fundamentale în informatică. Ele permit
modelarea și rezolvarea unor probleme complexe precum căutarea, optimizarea,
drumurile minime, conectivitatea și ordonarea dependențelor. Algoritmii
prezentați în acest capitol sunt esențiali pentru orice programator sau
informatician.

\end{document}
